As Large Language Models (LLMs) are increasingly integrated into critical decision-making processes, distinguishing between different modes of untruthfulness—specifically epistemic failure (hallucination) and strategic deception (intentional misreporting)—is paramount for safety and alignment. Current black-box lie detection methods have demonstrated significant promise in identifying falsehoods, yet it remains unclear whether they detect the act of deception or the falsity of the statement itself. In this work, we systematically investigate this distinction using the \datasetname dataset. By contrasting scenarios where \modelname either intentionally lies about known facts or inadvertently hallucinates unknown ones, we evaluate the specificity of elicitation-based lie detectors. Our results reveal a stark divergence: while intentional deception is highly detectable with a mean \liarscore of 6.13 ($p < 0.0001$), hallucinations remain largely "invisible" to these detectors, yielding scores (3.20) statistically indistinguishable from truthful controls (2.60). These findings suggest that current behavioral lie detectors primarily monitor the model's internal awareness of untruthfulness rather than external factual accuracy. We conclude that robust AI monitoring systems must decouple deception detection from epistemic validation to ensure comprehensive reliability.
